\section{Preliminaries}




\subsection{Reinforcement Learning with Verifiable Rewards for Vision-Language Models}

The reinforcement learning with verifiable rewards (RLVR) framework for vision-language models (VLMs) closely parallels that for large language models (LLMs), with the primary distinction being that RLVR for VLMs processes both an image and a text query as input to generate a textual chain-of-thought culminating in a verifiable answer prediction, whereas RLVR for LLMs operates solely on text queries.
Specifically, RLVR for VLMs aims to maximize the following objective:
\begin{align}
    &J(\pi) = \mathbb{E}_{(\boldsymbol{I}, \boldsymbol{q}, \boldsymbol{a}) \sim \mathcal{D}, \boldsymbol{o} \sim \pi(\cdot \mid \boldsymbol{I}, \boldsymbol{q})}[R(\boldsymbol{o}, \boldsymbol{a})],
    \\
    \text{where} & \quad  R(\boldsymbol{o}, \boldsymbol{a}) = \begin{cases} 1.0 & \text{if}\ \texttt{is\_equivalent}(\boldsymbol{o}, \boldsymbol{a}), \\ 0.0 & \text{otherwise.} \end{cases} \label{eq: reward calculation}
\end{align}
Here, $\boldsymbol{I}$, $\boldsymbol{q}$, and $\boldsymbol{a}$ denote the image, text query, and ground-truth answer, respectively, sampled from dataset $\mathcal{D}$, and $\boldsymbol{o}$ represents the response generated by policy $\pi$ conditioned on $\boldsymbol{I}$ and $\boldsymbol{q}$.




\subsection{Soft Adaptive Policy Optimization}

Soft Adaptive Policy Optimization (SAPO, \CiteParen{sapo}) is introduced to mitigate the potential instability and inefficiency caused by hard clipping in prior RLVR algorithms such as GSPO~\CiteParen{gspo} and GRPO~\CiteParen{shao2024deepseekmath}.
Concretely, SAPO substitutes hard clipping with a temperature-controlled soft gate and optimizes the following objective for VLMs:
\begin{equation}
\mathcal{J}(\theta)
=
\mathbb{E}_{(\boldsymbol{I}, \boldsymbol{q}, \boldsymbol{a}) \sim \mathcal{D}, \{\boldsymbol{o}_i\}_{i=1}^G \sim \pi_{\text{old}}(\cdot \mid \boldsymbol{I}, \boldsymbol{q})}
\left[
\frac{1}{G}
\sum_{i=1}^{G}
\frac{1}{|\boldsymbol{o}_i|}
\sum_{t=1}^{|\boldsymbol{o}_i|}
f_{i,t}\!\left(r_{i,t}(\theta)\right)
\hat{A}_{i,t}
\right],
\end{equation}
\begin{align}
\text{where} \quad &r_{i,t}(\theta)
=
\frac{
\pi_{\theta}\!\left(\boldsymbol{o}_{i,t} \mid \boldsymbol{I}, \boldsymbol{q}, \boldsymbol{o}_{i,<t}\right)
}{
\pi_{\theta_{\text{old}}}\!\left(\boldsymbol{o}_{i,t} \mid \boldsymbol{I}, \boldsymbol{q}, \boldsymbol{o}_{i,<t}\right)
},
\qquad
\hat{A}_{i,t}
=
\hat{A}_i
=
\frac{
R_i - \operatorname{mean}\!\left(\{R_j\}_{j=1}^{G}\right)
}{
\operatorname{std}\!\left(\{R_j\}_{j=1}^{G}\right)
},
\\
&f_{i,t}(x)
=
\sigma\!\left(\tau_{i,t}(x - 1)\right)
\cdot
\frac{4}{\tau_{i,t}},
\qquad
\tau_{i,t}
=
\begin{cases}
\tau_{\text{pos}}, & \text{if } \hat{A}_{i,t} > 0, \\
\tau_{\text{neg}}, & \text{otherwise}.
\end{cases}
\end{align}
Here, $i, j$ denote the sample indices in $\mathcal{D}$, and $t$ is the token index within a sequence. 
The parameters of the currently trained policy and the old rollout policy are denoted by $\theta$ and $\theta_{\text{old}}$, respectively. 
The temperatures for positive and negative tokens are $\tau_{\text{pos}}$ and $\tau_{\text{neg}}$, respectively. 
Moreover, $R_i$ is computed according to \cref{eq: reward calculation} for the $i$-th sample, and $\sigma(x) = \frac{1}{1 + e^{-x}}$ denotes the sigmoid function.
